\section{安全模型：从争议到可操作防线}

访谈后半段大量讨论安全问题，信息密度很高。一个关键点是：\textbf{OpenClaw 的风险不是传统 Web 漏洞的简单放大，而是“高权限 + 自治行为”带来的新型复合风险}。

\subsection{主要风险面}

\begin{warningbox}{OpenClaw 风险面的四个入口}
\begin{enumerate}
    \item \textbf{Inbound Prompt}：外部输入注入；
    \item \textbf{Tool Blast Radius}：工具调用越权；
    \item \textbf{Network Exposure}：外网访问与数据外泄；
    \item \textbf{Local Disk / Memory}：本地状态持久化污染。
\end{enumerate}
\end{warningbox}

\begin{importantbox}{访谈中提到的防护思路}
\begin{itemize}
    \item 提供安全审计脚本，先做配置检查；
    \item 默认提醒用户不要盲信低质量模型；
    \item 对高权限动作增加明确授权；
    \item 将日志和记忆写入可检查位置，避免“黑箱决策”。
\end{itemize}
\end{importantbox}

\subsection{“安全社区围攻”也是加速器}

Peter 的表述很现实：被大量安全研究者同时审视很痛苦，但也相当于获得了高强度免费审计。对早期项目，这是压力也是机会。

\begin{knowledgebox}{安全节奏建议}
对早期 agent 项目，建议采用“能力开放节奏 < 安全收敛节奏”原则：先限制 blast radius，再逐步放宽自治权限，而不是反过来。
\end{knowledgebox}

\begin{figure}[H]
\centering
\includegraphics[width=0.9\textwidth]{frames/frame-025932.jpg}
\caption{安全治理讨论：审计、暴露面与权限边界\protect\footnotemark}
\end{figure}
\footnotetext{视频画面时间区间：02:59:25--02:59:40。}

\subsection{本章小结}
OpenClaw 的安全议题说明，agent 产品化不是“做个更强模型”就能解决。真正难点在于\textbf{把能力映射为可治理权限}，并持续收敛风险面。

